\chapter{Navier-Stokes}
\label{chap:navier-stokes}

\section{Adaptive Grids}

Many systems exhibit a multi-scale nature, where phenomenon occur at a
hierarchy of different physical and temporal scales. This observation,
particularly in physical space, has been the motivation behind the development
of various \emph{adaptive mesh refinement} (AMR) schemes. The general idea
across different forms of AMR is to capture small-scale phenomenon with fine
meshes \emph{locally} and to keep solution meshes coarse where phenomenon occur
on a larger scale. Ideally, these adaptive meshes are then able to dynamically
adapt the grid (through refinement or coarsening) throughout the course of the
simulation as according to user-defined and problem-specific refinement
criteria.

For Cartesian grids, the oldest and arguably still most popular type of
approaches are broadly called \emph{block-structured} AMR. This body of methods
originates from a series of papers \cite{berger_adaptive_1984} and still
remains popular and the family of schemes is still widely used in modern
libraries (e.g. \href{https://amrex-codes.github.io/amrex/}{AMReX}). In this
approach, a coarse Cartesian grid is created with size \(h\), where regions
requiring refinement have finer rectangular grids \(h/2\) nested within the
coarse grid.

A second broad category of method applicable to cartesian grids is
\emph{tree-structured} AMR. These methods represent a Cartesian mesh based on a
tree data structure, where all possible cells or nodes of a mesh are
represented as the \emph{leaves} of a tree. The main functional difference of
these methods is the creation of locally-refined regions through a
``point-wise'' refinement the size of a single mesh cell, rather than through
defining larger rectangular regions as is required using block-structured AMR.
In theory, this allows more flexible and local refinement in regions of
interest with more arbitrary geometry, in comparison to defining regions of
interest only though blocks. The trade off is an arguably more complex
implementation and less inherent cache-locality. In addition, special treatment
is needed for the \emph{restriction} and \emph{prolongation} operators such
that schemes using the tree-structured mesh retain desired order near
refinement boundaries.

\section{Tree-structured Adaptive Meshes}

Tree-structured adaptive meshes may be understood as hierarchical data
structures for representing a computational domain by recursively subdividing
cells. In a single-tree grid, the \emph{root} cell typically coincides with the
entire physical domain. Refinement proceeds by replacing a cell with \(2^d\)
children, where \(d\) is the spatial dimension. Thus, in two dimensions the
resulting structure is a \emph{quadtree}, while in three dimensions it is an
\emph{octree}. Cells without children are called the \emph{leaves} of the tree
and form the active mesh on which the discrete equations are usually defined.
As with other abstract data structures, different implementations may realize
the same logical tree in different ways, with corresponding trade-offs in
memory layout, traversal efficiency, parallel scalability, and ease of neighbor
lookup.

Before tree-adaptive meshes became common in scientific computing, they were
studied first as general-purpose spatial data structures starting around the
late 1970s. Their recursive refinement makes them a compact representation of
spatially localized detail: uniform regions can be stored coarsely, while
geometric features, interfaces, or occupied regions can be represented at
higher resolution. Consequently, early applications appeared in computer
graphics, image and volume representation, solid modeling, spatial indexing,
neighbor queries, visibility algorithms, and ray-tracing acceleration. In these
settings, the leaves of the tree represented regions of space for storage,
search, or geometric reasoning, rather than finite-volume or finite-element
cells for solving partial differential equations. Practical application of the
data structure would only be realized in the early 2000s, such as in RAMSES and
Gerris.

\subsection{Data Structures for Adaptive-Tree Meshes}

One naive implementation of the quad- or octree datastructure is through a
struct or class representing each node of the tree, where each vertex contains
\(2^d\) pointers in memory to each child. This pointer then in effect
represents the directed edge to the children.

\paragraph{Fully Threaded Tree}
An extension of the naive implmentation is the \emph{full-threaded tree}
described by \textcite{khokhlov_fully_1998}. In addition to the struct or class
holding a pointer to the \(2^d\) children, we may also have each node contain
pointers to the parent, as well as each neighbor cell connected by a face or
edge. The obvious benefit of to a computational mesh can be seen if
implementing a finite-difference stencil, where frequent neighbor traversal is
needed, which does not come for free from simple pointer-arithmetic such as
when traversing a simple array.

Unless special care is taken, one of the potential downsides of these types of
trees is the memory requirement of storing pointers, as well as the
non-inherent memory locality between cells. Although the \(2^d\) children may
be heap-allocated and be adjacent in memory, leaves without common parents are
in no way guarenteed to be nearby in memory. This is in contrast to simple
arrays, where, since the entire grid is allocated at once on the heap, spatial
neighbors are local in memory also, and allow modern computers and compilers to
take advantage of CPU cache and vetorization.

\paragraph{Linear Trees}
Another approach to representing trees is the \emph{linear tree}. Rather than
representing the tree as simple structs or classes connected by pointers, we
may store each leaf of the tree in a linear array, where each element of the
array stores a spatial index encoded in a single binary integer. The entire
abstract tree, including neighbor queries, then may be recovered purely through
bit operations and array indicies.

These representations of trees have several benefits: One is a reduction in the
memory storage requirements, due to the fact that the edges and threads can be
recovered from a single integer. Second, is that when sorted by the Morton or
Z-order code, computations performed on the vertices (e.g. with a stencil),
provide some guarantee that cells close together in space will also be close
together in memory, which is good for avoiding {cache misses} and allow
efficient vectorization of operations. Another benefit of using this encoding,
is that it allows a straightforward method for domain-decomposition: Simply
divide the spatial domain by splitting the sorted array into \(n\) sub-arrays.
The current state-of-the-art in this regard may be found in
\textcite{burstedde_p4est_2011}.

The primary and still very acceptable disadvantages are two: The maximum depth
of your domain is now limited by the size of the integer used to represent the
spatial code. Second, and perhaps more importantly, insertion (refinement),
deletion (coarsening) and neighbour look-ups now become
\(\mathcal{O}(\log{(n)})\) operations (where \(n\) is the number of leaf
vertices): This is due to the fact one now must perform a binary search through
the sorted array of leaf vertices, where in pure pointer-based trees, these
operations are principally \(\mathcal{O}(1)\).

An example of such a scheme is illustrated in Figure \ref{fig:linear-quadtree}
where a sample 6-digit binary string is used to encode a representation of a
simple quadtree mesh.

\begin{figure}[H]
    \centering
    \def\svgwidth{\textwidth}
    \import{figure/vector}{linear-quadtree.pdf_tex}
    \caption{An example of a linear quadtree implementation where cells positions are given an \((i,j)\) index in an coordinate system with origin defined similarly to matrix row-column indexing. Once a \((i,j)\) coordinate is assigned to each leaf cell of the tree (visible in the mesh on the left), the coordinate, together with the level, is used to compute a Morton-order code, shown in hexadecimal format in the mesh and tree, and in binary in the sorted array. This code is formed by interleaving the \(j\) and \(i\) coordinate in binary as shown in the diagram on the right. Once a code is generated for each leaf, the array on the right is sorted by this code, which corresponds to the dotted line in red.}
    \label{fig:linear-quadtree}
\end{figure}

\paragraph{Basilisk's Data Structures}

Unlike linear quadtree implementations such as \textcite{burstedde_p4est_2011},
which focuses strongly on meshing and domain decomposion concerns, Basilisk
presents itself as complete solver-development framework, providing facilities
for both adaptive meshing as well as tooling to support the discrete numerical
schemes on these adaptive Cartesian meshes as well. Since these concerns are
more tightly coupled in Basilisk, design of the data structures in Basilisk
routines favor a more monolothic design that manages to address both mesh
topology and solver concerns together. As a result, the scheme is closer in
shape to the FTT of \textcite{khokhlov_fully_1998} rather than the linear
quadtree of \textcite{burstedde_p4est_2011}.

In Basilisk, rather than using the default memory allocator, where memory is
given at random addresses as provided by the operating system, an
\emph{mempool} allocator\footnote{\url{https://basilisk.fr/src/grid/mempool.h}}
is used to allocate raw memory at each level. The mempool requests large chunks
of memory at a time, up to \(2^{20}\) bytes at a time. When a cell is refined,
the \texttt{refine} function asks the mempool for an address to the first free
and sufficiently large segment of memory availible in the pool. Then, the new
cells are written into the free block of memory in that pool. Similarly, when
\texttt{coarsen} is called, it informs the mempool that the block is no longer
used and may be reused at a later time. Actual information written into the raw
memory block include a small \texttt{Cell} struct containing the Cartesian
index \((i,j,k)\) followed by a given number of \texttt{double} variables
needed to store all Eulerian fields required by any particular solver. The
advantage of this is that although the non-sibling neghbors of cells are not
arranged or sorted in memory by means of a space-filling curve, they are
nontheless likely to be found within the same block of memory and therefore not
frequently cause cache misses.

However, since they are not arranged by spatial ordering, a separate means of
neighbor finding is needed. Basilisk does this through a memory
index\footnote{\url{https://basilisk.fr/src/grid/memindex/virtual.h}}. At each
refinement level, Basilisk stores a \texttt{Layer} containing both the
\texttt{Mempool} which owns the cell data and a \texttt{Memindex} which maps
logical Cartesian indices to addresses in that pool. The \texttt{Memindex} is
therefore not a space-filling ordering of the cells, but a sparse
multidimensional array of pointers keyed by the integer coordinates \((i,j,k)\)
of the level. When children are created, Basilisk obtains a contiguous block
from the level's mempool and assigns each child coordinate to the corresponding
address in the index; when children are coarsened, these entries are cleared.
Neighbor access can then be implemented by adding an offset to the current
point's integer coordinates and dereferencing the corresponding entry, rather
than by searching a sorted Morton array, making the operation overall
\(\mathcal{O}(1)\). The virtual-memory implementation uses \texttt{mmap} to
allocate the pointer tables lazily and \texttt{madvise} to release pages when
they become empty, so the code can expose the full Cartesian index space of a
level, including ghost cells, without storing a dense array of cell data.

\subsection{Tree Grading and Interpolation Schemes}

In all adaptive mesh schemes, special treatement is needed at refinement
boundaries. Though it is factually true that both block-structured AMR and
tree-structured AMR can represent the exact same topologies, the primary
limitation that practically restricts either to a subset of ungraded meshes is
the interpolation schemes used at refinement boundaries. Typically, the
block-structured AMR methods impose minimum patch sizes because they allow for
much simpler interpolation schemes at the cost of stronger restrictions on the
types of mesh topologies allowed. Similarly, in tree-structured AMR, it is
typically imposed that cells differ at most 1 level across a shared face, edge,
or vertex. Typically, this is stated as a "2:1" grading. This is entirely done,
as in the patch-based AMR, to ensure that ghost cells may be interpolated at
refinement boundaries, usually in a way that preserves local or global order
for a given scheme. A good summary of this particular issue in adaptive meshing
may be found in \textcite{batty_cell-centred_2017} and the particular
interpolation scheme still used by Basilisk may be found in
\textcite{popinet_gerris_2003}.

\subsection{Domain Decomposition}

\section{Centered Projection Solver}
\label{sec:centered-projection-solver}

The centered Navier--Stokes module in Basilisk is a finite-volume projection
method for the variable-density Navier--Stokes
equations.\footnote{\url{https://basilisk.fr/src/navier-stokes/centered.h}} The
primary unknowns are the cell-centered velocity \(\vec u\) and pressure \(p\).
The solver also maintains a face-centered velocity \(\vec u_f\), used for
fluxes and for the pressure projection, and a cell-centered vector \(\vec G\)
which stores the combined pressure-gradient and acceleration contribution. In
continuum form, the equations approximated by the solver are
\begin{align}
    \pdv{\vec u}{t}
    + \nabla \cdot \left(\vec u \otimes \vec u\right)
                        & =
    \frac{1}{\rho_f}
    \left[-\nabla p + \nabla \cdot \left(2\mu_f \mat D\right)\right]
    + \vec a_f,
    \label{eq:basilisk-centered-momentum} \\
    \nabla \cdot \vec u & = 0,
    \label{eq:basilisk-centered-incompressibility}
\end{align}
where
\begin{equation}
    \mat D = \frac{1}{2}\left[\nabla \vec u
        + \left(\nabla \vec u\right)^\transp\right]
\end{equation}
is the rate-of-strain tensor. In the discrete solver, the specific volume
\(\alpha_f = 1/\rho_f\), dynamic viscosity \(\mu_f\), and acceleration \(\vec a_f\)
are stored on cell faces. If the user does not provide these fields, the
default state is a unit-density inviscid flow with zero acceleration. The
implementation follows the projection-method structure used in Gerris
\cite{popinet_gerris_2003} and Basilisk's adaptive finite-volume solvers
\cite{popinet_quadtree-adaptive_2015}.

The pressure and acceleration terms are discretized consistently on faces. On a
regular Cartesian grid, the face value of the combined source in the
\(x\)-direction is
\begin{equation}
    \left(G_f\right)^x_{i+1/2}
    =
    a^x_{i+1/2}
    - \alpha^x_{i+1/2}\frac{p_{i+1}-p_i}{\Delta x}.
    \label{eq:basilisk-face-gradient}
\end{equation}
The cell-centered field used to update \(\vec u\) is then obtained by
averaging the adjacent face values,
\begin{equation}
    G^x_i =
    \frac{1}{2}
    \left[\left(G_f\right)^x_{i-1/2}
    + \left(G_f\right)^x_{i+1/2}\right],
    \label{eq:basilisk-centered-gradient}
\end{equation}
with analogous expressions in the other coordinate directions. This
face-based construction is part of Basilisk's source-term interface: modules
for surface tension, reduced gravity, buoyancy, and user-defined body forces
typically add their contributions through the face acceleration field
\(\vec a_f\). Combining \(\vec a_f\) with the pressure gradient on the same face
stencil allows discrete equilibria, such as Laplace-pressure balance or
hydrostatic balance, to be represented without introducing a spurious velocity
field. At solid or symmetry boundaries, the default pressure condition is
therefore a Neumann condition chosen so that the normal pressure gradient
balances the normal acceleration.

The projection operator acts on a provisional face velocity \(\vec w_f^\star\).
For a projection over a time interval \(\Delta \tau\), the pressure-like scalar
\(\phi\) satisfies
\begin{equation}
    \nabla_h \cdot \left(\alpha_f \nabla_h \phi\right)
    =
    \frac{1}{\Delta \tau}\nabla_h \cdot \vec w_f^\star,
    \label{eq:basilisk-projection-poisson}
\end{equation}
and the projected face velocity is
\begin{equation}
    \vec w_f
    =
    \vec w_f^\star
    - \Delta \tau\,\alpha_f\nabla_h\phi .
    \label{eq:basilisk-projection-update}
\end{equation}
This enforces \(\nabla_h\cdot\vec w_f=0\) up to the tolerance of the
multigrid Poisson solve. The method is therefore an approximate projection:
the face velocity is projected directly, while the cell-centered velocity is
corrected using the averaged field \(\vec G\) from
\eqref{eq:basilisk-centered-gradient}.

The half-step projection in the advection stage supplies a divergence-free
advecting velocity for the Bell--Colella--Glaz fluxes. The final projection
then computes the pressure at \(t^{n+1}\), updates the face velocity used by
the next time step, and applies the corresponding centered pressure and
acceleration correction to \(\vec u\). When the \texttt{stokes} flag is set,
the explicit advection stage is skipped and the method reduces to the viscous,
pressure, and acceleration updates. On tree-adaptive grids, Basilisk uses
restriction and prolongation operators chosen to preserve the solenoidal
character of the face velocity across refinement boundaries.

In Algorithm~\ref{alg:basilisk-centered-solver}, \(\mathcal P_{\Delta\tau}\)
denotes the projection defined by
\eqref{eq:basilisk-projection-poisson}--\eqref{eq:basilisk-projection-update},
\(\mathcal V_{\Delta t}\) denotes the implicit viscous solve, \(I_f\) is the
cell-to-face interpolation of the centered velocity, and \(\mathcal G(p,\vec
a_f)\) is the centered pressure/acceleration operator in
\eqref{eq:basilisk-face-gradient}--\eqref{eq:basilisk-centered-gradient}.

\begin{algorithm}[H]
    \caption{Pseudocode for one centered Navier--Stokes time step}
    \label{alg:basilisk-centered-solver}
    \DontPrintSemicolon
    \KwIn{\(\vec u^n\), \(p^n\), \(\vec G^n\), \(\vec u_f^n\), and
        material/source fields \(\alpha_f\), \(\rho_f\), \(\mu_f\), \(\vec a_f\).}
    \KwOut{\(\vec u^{n+1}\), \(p^{n+1}\), \(\vec G^{n+1}\), and
    \(\vec u_f^{n+1}\) with \(\nabla_h\cdot\vec u_f^{n+1}=0\).}
    \(\Delta t \leftarrow \textsc{CFL}(\vec u_f^n)\)\;
    \(\textsc{AdvanceTracers}(\Delta t/2)\)\;
    \((\alpha_f,\rho_f,\mu_f) \leftarrow \textsc{Properties}()\)\;
    \(\vec u^\star \leftarrow \vec u^n\)\;
    \If{\(\texttt{stokes}=\texttt{false}\)}{
    \(\widehat{\vec u}_f^{\,n+1/2}
    \leftarrow \textsc{BCGPredict}(\vec u^n,\vec G^n,\Delta t)\)\;
    \((\vec u_f^{\,n+1/2},p_f)
    \leftarrow
    \mathcal P_{\Delta t/2}(\widehat{\vec u}_f^{\,n+1/2};\alpha_f)\)\;
    \(\vec u^\star
    \leftarrow
    \textsc{BCGAdvect}(\vec u^n,\vec u_f^{\,n+1/2},\vec G^n,\Delta t)\)\;
    }
    \If{\(\mu_f \ne 0\)}{
        \(\vec u^\star
        \leftarrow
        \mathcal V_{\Delta t}(\vec u^\star+\Delta t\,\vec G^n;\rho_f,\mu_f)
        -\Delta t\,\vec G^n\)\;
    }
    \(\vec a_f \leftarrow \textsc{AccelerationSources}()\)\;
    \(\widehat{\vec u}_f^{\,n+1}
    \leftarrow I_f(\vec u^\star)+\Delta t\,\vec a_f\)\;
    \((\vec u_f^{\,n+1},p^{n+1})
    \leftarrow
    \mathcal P_{\Delta t}(\widehat{\vec u}_f^{\,n+1};\alpha_f)\)\;
    \(\vec G^{n+1} \leftarrow \mathcal G(p^{n+1},\vec a_f)\)\;
    \(\vec u^{n+1} \leftarrow \vec u^\star+\Delta t\,\vec G^{n+1}\)\;
\end{algorithm}
